\section{Keeping track of what depends on what}\label{ch01:sec:dependence}
In the proof of Proposition~\ref{thm:odd-sum}, $a$ and $b$ can be any odd integers, but the witnesses are not chosen freely. Solving $a=2k+1$ gives $k=(a-1)/2$, and likewise $\ell=(b-1)/2$, so each witness is determined by its odd integer. If $a$ changes from $7$ to $9$ while $b$ stays the same, then $k$ changes from $3$ to $4$, but $\ell$ does not. The witness $k+\ell+1$ for the sum is built in turn from $k$ and $\ell$.

Suppose an assistant writes both odd integers as $2k+1$, using the same $k$ for each, and computes $(2k+1)+(2k+1)=2(2k+1)$, which is certainly even. The algebra is correct, but the argument covers only the case in which the two odd integers are equal, as in $5+5$. It says nothing about $3+5$, because no single integer $k$ makes both $3=2k+1$ and $5=2k+1$ true. So ask of any argument: ``Which quantities can be anything, and which are determined by the others?'' Chapter~\ref{ch:logic} makes this question precise with the phrases ``for every'' and ``there exists.''

Separate witnesses matter just as much when we show that the sum of three odd integers is odd. Let $a$, $b$, and $c$ be odd integers. By the definition of odd, there are integers $k$, $\ell$, and $m$, possibly different, with $a=2k+1$, $b=2\ell+1$, and $c=2m+1$. Adding gives
\[
a+b+c=(2k+1)+(2\ell+1)+(2m+1)=2k+2\ell+2m+3.
\]
This may ``look odd,'' but the definition asks for the form \emph{twice an integer, plus one}. So split $3$ as $2+1$ and factor $2$ out of the terms that contain it:
\[
2k+2\ell+2m+3=2k+2\ell+2m+2+1=2(k+\ell+m+1)+1.
\]
The number $k+\ell+m+1$ is an integer, so $a+b+c$ is odd by Definition~\ref{def:parity}.
